\documentclass[11pt,a4paper]{article}

% --- Idioma e codificação ---
\usepackage[utf8]{inputenc}
\usepackage[T1]{fontenc}
\usepackage[brazil]{babel}

% --- Layout ---
\usepackage[margin=2.5cm]{geometry}
\usepackage{parskip}
\usepackage{titlesec}
\usepackage{fancyhdr}

% --- Tabelas, cores e código ---
\usepackage{booktabs}
\usepackage{longtable}
\usepackage{array}
\usepackage{xcolor}
\usepackage{listings}

% --- Links e metadados ---
\usepackage{hyperref}
\hypersetup{
    colorlinks=true,
    linkcolor=blue!50!black,
    urlcolor=blue!50!black,
    citecolor=blue!50!black,
    pdftitle={Sucuri --- Relatório de análises: Despesas com Ensino Superior},
    pdfauthor={Projeto Sucuri}
}

\definecolor{codebg}{RGB}{245,245,245}
\lstset{
    backgroundcolor=\color{codebg},
    basicstyle=\ttfamily\small,
    breaklines=true,
    frame=single,
    showstringspaces=false,
    columns=fullflexible,
}

\pagestyle{fancy}
\fancyhf{}
\fancyhead[L]{Sucuri --- Detecção de anomalias em despesas públicas}
\fancyhead[R]{\thepage}
\renewcommand{\headrulewidth}{0.4pt}

\newcommand{\aceite}{\textcolor{green!45!black}{\textbf{OK}}}

\title{\textbf{Sucuri}\\[0.3em]
\large Detecção de anomalias em despesas federais com Ensino Superior\\
\large Relatório consolidado de análises, medidas e interpretações}
\author{Portal da Transparência (gov.br) --- Ministério da Educação}
\date{Atualizado em 16 de julho de 2026}

\begin{document}

\maketitle

\begin{abstract}
\noindent Este documento é a versão em \LaTeX{} de \texttt{relatorios/RELATORIO.md}
e deve conter, ao final do projeto, todas as análises, medidas estatísticas e
interpretações produzidas ao longo das fases do \texttt{ROADMAP.md}. A versão
atual documenta exclusivamente a \textbf{Fase 0 --- Infraestrutura do
projeto}, que prepara a base de código para as análises de dados das fases
seguintes; nenhuma conclusão sobre despesas é apresentada aqui ainda.
\end{abstract}

\tableofcontents
\vspace{1em}

\section{Objetivo desta fase}

A Fase 0 do ROADMAP não produz análises sobre as despesas em si --- seu
objetivo é transformar o projeto de um único script
(\texttt{coletar\_despesas.py}) em uma base de código reprodutível, testável
e reutilizável pelas fases seguintes (qualidade de dados, deflação, detecção
de anomalias). Esta seção documenta \textbf{o que foi feito}; interpretações
sobre os dados de despesas começam na Fase~1/2 e serão anexadas como novas
seções deste mesmo documento.

\section{Tarefa 0.1 --- Empacotamento e ambiente}

\subsection{Estrutura criada}

\begin{lstlisting}[language=bash]
pyproject.toml            # projeto gerenciado por uv (PEP 621 + hatchling)
uv.lock                   # lockfile de dependencias
src/sucuri/
  __init__.py
  api.py                  # cliente da API do Portal da Transparencia
  features.py             # engenharia de variaveis e construcao dos paineis A/B
  persistencia.py         # salvamento de CSV/Parquet/JSON bruto e dicionario
  utils.py                # brl_para_float, razao_segura, classificar_instituicao
analises/                 # reservado para as fases 1-2 (vazio nesta fase)
relatorios/
  figuras/                # reservado para as fases 1-2 (vazio nesta fase)
  latex/                  # este relatorio em LaTeX/PDF
tests/
  __init__.py
  test_utils.py
  test_features.py
coletar_despesas.py       # CLI fino: apenas argparse + orquestracao
\end{lstlisting}

\subsection{Migração de código}

Toda a lógica reutilizável de \texttt{coletar\_despesas.py} foi movida para o
pacote \texttt{src/sucuri/}, preservando exatamente o comportamento original
(mesmos parâmetros de API, mesmas fórmulas, mesmo formato de saída).
A Tabela~\ref{tab:migracao} resume o mapeamento.

\begin{table}[h]
\centering
\small
\begin{tabular}{@{}p{4.2cm}p{4.2cm}p{4.6cm}@{}}
\toprule
\textbf{Função/constante} & \textbf{Origem (script único)} & \textbf{Destino} \\
\midrule
\texttt{brl\_para\_float}, \texttt{\_razao\_segura}, \texttt{classificar\_instituicao} &
topo do script &
\texttt{sucuri/utils.py} (\texttt{razao\_segura} tornada pública) \\
\addlinespace
\texttt{carregar\_chave\_api}, \texttt{criar\_sessao}, \texttt{requisitar}, \texttt{coletar\_paginado}, \texttt{coletar\_intervalo} &
seção ``Coleta genérica'' &
\texttt{sucuri/api.py} \\
\addlinespace
\texttt{construir\_df\_funcional}, \texttt{construir\_df\_instituicao}, \texttt{\_indicadores\_execucao}, \texttt{\_features\_serie\_temporal}, \texttt{\_consolidar\_flags} &
seções (A)/(B) e ``Engenharia de variáveis'' &
\texttt{sucuri/features.py} (funções tornadas públicas) \\
\addlinespace
\texttt{salvar\_dataset}, \texttt{escrever\_dicionario}, \texttt{DICIONARIO}, \texttt{resumir} &
seção ``Persistência'' &
\texttt{sucuri/persistencia.py} \\
\bottomrule
\end{tabular}
\caption{Mapeamento das funções migradas do script único para o pacote \texttt{sucuri}.}
\label{tab:migracao}
\end{table}

O arquivo \texttt{coletar\_despesas.py} (antes
\texttt{coletar\_despesas\_ensino\_superior.py}) ficou reduzido a
aproximadamente 125 linhas: apenas \texttt{argparse}, orquestração das
coletas (A) e (B), e chamadas às funções do pacote. Nenhuma fórmula ou regra
de negócio foi alterada durante a migração --- apenas reorganizada em
módulos.

\subsection{Dependências}

Gerenciadas por \texttt{uv}: \texttt{pandas}, \texttt{pyarrow},
\texttt{requests}, \texttt{python-dotenv}, \texttt{scikit-learn},
\texttt{scipy}, \texttt{matplotlib} (produção); \texttt{ruff},
\texttt{pytest} (grupo \texttt{dev}). O comando \texttt{uv sync --group dev}
cria o ambiente virtual e resolve o lockfile em poucos segundos.

\subsection{Verificação dos critérios de aceite}

\begin{table}[h]
\centering
\begin{tabular}{@{}p{9.5cm}c@{}}
\toprule
\textbf{Critério (ROADMAP 0.1)} & \textbf{Resultado} \\
\midrule
\texttt{uv run python coletar\_despesas.py --help} funciona & \aceite \\
\texttt{uv run pytest} roda (mesmo com 0 testes) & \aceite \\
\texttt{uv run ruff check .} sem erros & \aceite \\
\bottomrule
\end{tabular}
\end{table}

\section{Tarefa 0.2 --- Testes unitários das funções críticas}

\subsection{Cobertura}

Foram escritos 37 testes em dois arquivos, cobrindo exatamente as funções
indicadas no ROADMAP.

\paragraph{\texttt{tests/test\_utils.py} (23 testes)}
\begin{itemize}
    \item \texttt{brl\_para\_float}: formato brasileiro com separador de
    milhar (\texttt{1.059.473.395,24}), sem separador de milhar,
    \texttt{None}, string vazia, string só com espaços, valor já numérico
    (\texttt{int}/\texttt{float}), entrada inválida (\texttt{"abc"}
    $\rightarrow$ \texttt{NaN}), valor negativo.
    \item \texttt{razao\_segura}: divisão normal e \textbf{denominador zero
    $\rightarrow$ \texttt{NaN}} (em vez de erro ou \texttt{inf}).
    \item \texttt{classificar\_instituicao}: um caso por categoria
    (Universidade Federal, Hospitalar/EBSERH, Instituto Federal, CEFET,
    CAPES, Fundo FNDE/FIES, Educação Básica, Outros/Administração), mais
    \texttt{None} e string vazia.
\end{itemize}

\paragraph{\texttt{tests/test\_features.py} (14 testes)}
\begin{itemize}
    \item \texttt{indicadores\_execucao}: taxas e saldos em caso normal;
    \textbf{\texttt{empenhado == 0} $\rightarrow$ taxas \texttt{NaN} sem
    exceção}; disparo isolado de cada flag de incoerência (\texttt{pago >
    empenhado}, \texttt{liquidado > empenhado}, valor negativo); ausência
    de falso positivo em valores coerentes.
    \item \texttt{features\_serie\_temporal}: variação anual ano a ano;
    disparo de \texttt{flag\_salto\_anual} acima de 100\%; \textbf{série
    com desvio-padrão zero (valores constantes) não gera erro de divisão
    por zero} e produz \texttt{zscore} \texttt{NaN} em vez de indefinido;
    série com uma única observação; independência entre séries diferentes
    agrupadas no mesmo DataFrame.
    \item \texttt{consolidar\_flags}: nenhuma flag ativa, uma flag ativa,
    flag extra (\texttt{flag\_atipico\_entre\_pares}) considerada, valores
    nulos tratados como \texttt{False} na consolidação.
\end{itemize}

Os dois casos de borda explicitamente pedidos no ROADMAP --- \textbf{divisão
por zero} e \textbf{série com desvio-padrão zero} --- estão cobertos e
passam sem exceções, confirmando que as fórmulas originais
(\texttt{.replace(0, pd.NA)}) já protegiam contra isso; os testes tornam
essa garantia explícita e verificável automaticamente.

\subsection{Resultado da execução}

\begin{lstlisting}[language=bash]
$ uv run pytest -v
============================== 37 passed in 0.67s ==============================

$ uv run ruff check .
All checks passed!
\end{lstlisting}

\textbf{Aceite (ROADMAP 0.2):} \aceite~\texttt{uv run pytest} verde; as
funções citadas no ROADMAP (\texttt{brl\_para\_float},
\texttt{classificar\_instituicao}, \texttt{\_indicadores\_execucao},
\texttt{\_features\_serie\_temporal}, \texttt{\_consolidar\_flags}) têm
cobertura direta.

\section{Validação de não regressão}

Após a migração, foi confirmado por busca no repositório que não restou
nenhuma referência às funções privadas antigas
(\texttt{\_indicadores\_execucao}, \texttt{\_features\_serie\_temporal},
\texttt{\_consolidar\_flags}, \texttt{\_razao\_segura}) fora do pacote
\texttt{sucuri}, e que os módulos migrados importam e executam corretamente.
Os dados já coletados em \texttt{dados/} (dois conjuntos, 2014--2026) não
foram alterados nesta fase --- nenhuma recoleta foi executada.

\section{Ressalvas e decisões registradas}

\begin{itemize}
    \item As funções antes privadas (\texttt{\_indicadores\_execucao} etc.)
    foram tornadas públicas (sem \texttt{\_}) ao mover para
    \texttt{sucuri/features.py}, pois passam a ser a API pública do pacote
    reutilizada pelas fases seguintes. Isso é uma mudança de visibilidade,
    não de comportamento.
    \item \texttt{analises/} e \texttt{relatorios/figuras/} foram criados
    vazios nesta fase --- serão populados a partir da Fase 1/2 do ROADMAP.
    \item Nenhum dado foi recoletado; a chave \texttt{GOVBR\_API\_KEY} não
    foi lida nem impressa durante esta fase (apenas os testes unitários,
    que não fazem chamadas de rede, foram executados).
\end{itemize}

\section{Estado do ROADMAP}

Tarefas 0.1 e 0.2 marcadas como concluídas em \texttt{ROADMAP.md}. Próxima
tarefa pendente: \textbf{1.1 Relatório de qualidade dos dados}
(\texttt{analises/01\_qualidade.py} $\rightarrow$
\texttt{relatorios/01\_qualidade.md}), que passa a ser a primeira seção
deste relatório a conter interpretação de dados reais de despesas.

\section*{Instrução de manutenção deste relatório}
\addcontentsline{toc}{section}{Instrução de manutenção deste relatório}

A partir da Fase 1, toda nova tarefa do ROADMAP que gerar análise, medida
estatística ou modelo deve acrescentar uma seção a este documento (e à
versão Markdown em \texttt{relatorios/RELATORIO.md}) com: o que foi medido,
o método usado, o resultado numérico e sua interpretação --- nunca apenas o
código ou o caminho do arquivo gerado. Ver \texttt{CLAUDE.md}, seção
``Registro de análises e relatórios''.

\end{document}
